%  LaTeX support: latex@mdpi.com 
%  For support, please attach all files needed for compiling as well as the log file, and specify your operating system, LaTeX version, and LaTeX editor.

%=================================================================
\documentclass[journal,article,submit,pdftex,moreauthors,mathematics]{Definitions/mdpi} 
%\documentclass[preprints,article,submit,pdftex,moreauthors]{Definitions/mdpi} 
% For posting an early version of this manuscript as a preprint, you may use "preprints" as the journal. Changing "submit" to "accept" before posting will remove line numbers.

% Below journals will use APA reference format:
% admsci, aieduc, behavsci, businesses, econometrics, economies, education, ejihpe, famsci, games, humans, ijcs, ijfs, jintelligence, journalmedia, jrfm, jsam, languages, peacestud, psycholint, publications, tourismhosp, youth

% Below journals will use Chicago reference format:
% arts, genealogy, histories, humanities, laws, literature, religions, risks, socsci

%--------------------
% Class Options:
%--------------------
%----------
% journal
%----------
% Choose between the following MDPI journals:
% accountaudit, acoustics, actuators, addictions, adhesives, admsci, adolescents, aerobiology, aerospace, agriculture, agriengineering, agrochemicals, agronomy, ai, aichem, aieduc, aieng, aimater, aimed, aipa, air, aisens, algorithms, allergies, alloys, amh, analog, analytica, analytics, anatomia, anesthres, animals, antibiotics, antibodies, antioxidants, applbiosci, appliedchem, appliedmath, appliedphys, applmech, applmicrobiol, applnano, applsci, aquacj, architecture, arm, arthropoda, arts, asc, asi, astronautics, astronomy, atmosphere, atoms, audiolres, automation, axioms, bacteria, batteries, bdcc, behavsci, beverages, biochem, bioengineering, biologics, biology, biomass, biomechanics, biomed, biomedicines, biomedinformatics, biomimetics, biomolecules, biophysica, bioresourbioprod, biosensors, biosphere, biotech, birds, blockchains, bloods, blsf, brainsci, breath, buildings, businesses, cancers, carbon, cardio, cardiogenetics, cardiovascmed, catalysts, cells, ceramics, challenges, chemengineering, chemistry, chemosensors, chemproc, children, chips, cimb, civileng, cleantechnol, climate, clinbioenerg, clinpract, clockssleep, cmd, cmtr, coasts, coatings, colloids, colorants, commodities, complexities, complications, compounds, computation, computers, condensedmatter, conservation, constrmater, cosmetics, covid, crops, cryo, cryptography, crystals, csmf, ctn, culture, curroncol, cyber, dairy, data, ddc, dentistry, dermato, dermatopathology, designs, devices, dhi, diabetology, diagnostics, dietetics, digital, disabilities, diseases, diversity, dna, drones, dynamics, earth, ebj, ecm, ecologies, econometrics, economies, edm, education, eesp, ejihpe, electricity, electrochem, electronicmat, electronics, encyclopedia, endocrines, energies, eng, engproc, entomology, entropy, environments, environremediat, epidemiologia, epigenomes, esa, est, famsci, fermentation, fibers, fintech, fire, fishes, fluids, foods, forecasting, forensicsci, forests, fossstud, foundations, fractalfract, fuels, future, futureinternet, futurepharmacol, futurephys, futuretransp, galaxies, games, gases, gastroent, gastrointestdisord, gastronomy, gels, genealogy, genes, geographies, geohazards, geomatics, geometry, geosciences, geotechnics, geriatrics, germs, glacies, grasses, green, greenhealth, gucdd, hardware, hazardousmatters, healthcare, hearts, hemato, hematolrep, hep, heritage, higheredu, highthroughput, histories, horticulturae, hospitals, humanities, humans, hydrobiology, hydrogen, hydrology, hydropower, hygiene, idr, iic, ijcs, ijem, ijerph, ijfs, ijgi, ijmd, ijms, ijns, ijom, ijpb, ijt, ijtm, ijtpp, ime, immuno, informatics, information, infrastructures, inorganics, insects, instruments, inventions, iot, j, jaestheticmed, jal, jcdd, jcm, jcp, jcrm, jcs, jcto, jdad, jdb, jdream, jemr, jeta, jfb, jfmk, jgbg, jgg, jimaging, jintelligence, jlpea, jmahp, jmmp, jmms, jmp, jmse, jne, jnt, jof, joi, joitmc, joma, jop, joptm, jor, journalmedia, jox, jpbi, jphytomed, jpm, jrfm, jsam, jsan, jtaer, jvd, jzbg, kidneydial, kinasesphosphatases, knowledge, labmed, laboratories, lae, land, languages, laws, life, lights, limnolrev, lipidology, liquids, literature, livers, logics, logistics, lubricants, lymphatics, machines, macromol, magnetism, magnetochemistry, make, marinedrugs, materials, materproc, mathematics, mca, measurements, medicina, medicines, medsci, membranes, merits, metabolites, metals, meteorology, methane, metrics, metrology, micro, microarrays, microbiolres, microelectronics, micromachines, microorganisms, microplastics, microwave, minerals, mining, mmphys, modelling, molbank, molecules, mps, msf, mti, multimedia, muscles, nanoenergyadv, nanomanufacturing, nanomaterials, ncrna, ndt, network, neuroglia, neuroimaging, neurolint, neurosci, nitrogen, notspecified, nursrep, nutraceuticals, nutrients, obesities, occuphealth, oceans, ohbm, onco, optics, oral, organics, organoids, osteology, oxygen, pandemics, parasites, parasitologia, particles, pathogens, pathophysiology, peacestud, pediatrrep, pets, pharmaceuticals, pharmaceutics, pharmacoepidemiology, pharmacy, philosophies, photochem, photonics, phycology, physchem, physics, physiologia, plants, plasma, platforms, pollutants, polymers, polysaccharides, populations, poultry, powders, precipitation, precisoncol, preprints, proceedings, processes, prosthesis, proteomes, psf, psychiatryint, psychoactives, psycholint, publications, purification, quantumrep, quaternary, qubs, radiation, rdt, reactions, realestate, receptors, recycling, regeneration, religions, remotesensing, reports, reprodmed, resources, rheumato, risks, rjpm, robotics, rsee, ruminants, safety, sci, scipharm, sclerosis, seeds, sensors, separations, sexes, shi, signals, sinusitis, siuj, skins, smartcities, sna, societies, socsci, software, soilsystems, solar, solids, spectroscj, sports, standards, stats, std, stratsediment, stresses, surfaces, surgeries, suschem, sustainability, symmetry, synbio, systems, tae, targets, taxonomy, technologies, telecom, test, textiles, thalassrep, therapeutics, thermo, timespace, tomography, tourismhosp, toxics, toxins, tph, transplantology, transportation, traumacare, traumas, tri, tropicalmed, universe, urbansci, uro, vaccines, vehicles, venereology, vetsci, vibration, virtualworlds, viruses, vision, waste, water, welding, wem, wevj, wild, wind, women, world, youth, zoonoticdis

%---------
% article
%---------
% The default type of manuscript is "article", but can be replaced by: 
% abstract, addendum, article, benchmark, book, bookreview, briefcommunication, briefreport, casereport, changes, clinicopathologicalchallenge, comment, commentary, communication, conceptpaper, conferenceproceedings, correction, conferencereport, creative, datadescriptor, discussion, entry, expressionofconcern, extendedabstract, editorial, essay, erratum, fieldguide, hypothesis, interestingimages, letter, meetingreport, monograph, newbookreceived, obituary, opinion, proceedingpaper, projectreport, reply, retraction, review, perspective, protocol, shortnote, studyprotocol, supfile, systematicreview, technicalnote, viewpoint, guidelines, registeredreport, tutorial,  giantsinurology, urologyaroundtheworld
% supfile = supplementary materials

%----------
% submit
%----------
% The class option "submit" will be changed to "accept" by the Editorial Office when the paper is accepted. This will only make changes to the frontpage (e.g., the logo of the journal will get visible), the headings, and the copyright information. Also, line numbering will be removed. Journal info and pagination for accepted papers will also be assigned by the Editorial Office.

%------------------
% moreauthors
%------------------
% If there is only one author the class option oneauthor should be used. Otherwise use the class option moreauthors.

%---------
% pdftex
%---------
% The option pdftex is for use with pdfLaTeX. Remove "pdftex" for (1) compiling with LaTeX & dvi2pdf (if eps figures are used) or for (2) compiling with XeLaTeX.

%=================================================================
% MDPI internal commands - do not modify
\firstpage{1} 
\makeatletter 
\setcounter{page}{\@firstpage} 
\makeatother
\pubvolume{1}
\issuenum{1}
\articlenumber{0}
\pubyear{2026}
\copyrightyear{2026}
%\externaleditor{Firstname Lastname} % More than 1 editor, please add `` and '' before the last editor name
\datereceived{ } 
\daterevised{ } % Comment out if no revised date
\dateaccepted{ } 
\datepublished{ } 
%\datecorrected{} % For corrected papers: "Corrected: XXX" date in the original paper.
%\dateretracted{} % For retracted papers: "Retracted: XXX" date in the original paper.
%\doinum{} % Used for some special journals, like molbank
%\pdfoutput=1 % Uncommented for upload to arXiv.org
%\CorrStatement{yes}  % For updates
%\longauthorlist{yes} % For many authors that exceed the left citation part
%\IsAssociation{yes} % For association journals

%=================================================================
% Add packages and commands here. The following packages are loaded in our class file: fontenc, inputenc, calc, indentfirst, fancyhdr, graphicx, epstopdf, lastpage, ifthen, float, amsmath, amssymb, lineno, setspace, enumitem, mathpazo, booktabs, titlesec, etoolbox, tabto, xcolor, colortbl, soul, multirow, microtype, tikz, totcount, changepage, attrib, upgreek, array, tabularx, pbox, ragged2e, tocloft, marginnote, marginfix, enotez, amsthm, natbib, hyperref, cleveref, scrextend, url, geometry, newfloat, caption, draftwatermark, seqsplit
% cleveref: load \crefname definitions after \begin{document}

%=================================================================
% Please use the following mathematics environments: Theorem, Lemma, Corollary, Proposition, Characterization, Property, Problem, Example, ExamplesandDefinitions, Hypothesis, Remark, Definition, Notation, Assumption
%% For proofs, please use the proof environment (the amsthm package is loaded by the MDPI class).

%=================================================================
% Full title of the paper (Capitalized)
\Title{CADEMAS-ML: A Software Tool for Context-Aware Cooperative Automated Decision-Making}

% Author Orchid ID: enter ID or remove command
\newcommand{\orcidauthorA}{0000-0003-3267-6753} % Add \orcidA{} behind the author's name
\newcommand{\orcidauthorB}{0009-0002-0183-7905} % Add \orcidB{} behind the author's name
\newcommand{\orcidauthorC}{0000-0002-7653-1452} % Add \orcidB{} behind the author's name

% Authors, for the paper (add full first names)
\Author{Pavel Novoa-Hernández$^{1,*}$\orcidA{}, Mariia Godz$^{2}$\orcidA{} and David A. Pelta$^{2}$\orcidC{}}

%\longauthorlist{yes}

% MDPI internal command: Authors, for metadata in PDF
\AuthorNames{Pavel Novoa-Hernández, Mariia Godz and David A. Pelta}

%\longauthorlist{yes}

% Affiliations / Addresses (Add [1] after \address if there is only one affiliation.)
\address{%
$^{1}$ \quad Dept. Computer Engineering and Systems, Universidad de La Laguna; pnovoahe@ull.edu.es\\
$^{2}$ \quad Dept. Computer Science and Artificial Intelligence, Universidad de Granada; \{mariiagodz,dpelta\}@ugr.es}

% Contact information of the corresponding author
\corres{Correspondence:pnovoahe@ull.edu.es}

% Current address and/or shared authorship
%\firstnote{Current address: Affiliation.}  
% Current address should not be the same as any items in the Affiliation section.

%\secondnote{These authors contributed equally to this work.}
% The commands \thirdnote{} till \eighthnote{} are available for further notes.

%\simplesumm{} % Simple summary

%\conference{} % An extended version of a conference paper

% Abstract (Do not insert blank lines, i.e. \\) 
\abstract{Automated decision-making systems are increasingly used in prediction and classification tasks. However, their outputs are often interpreted without considering the organizational, operational, or strategic context in which decisions are made. This paper presents CADEMAS-ML, a software tool for the design and analysis of cooperative and context-aware automated decision-making systems. The proposed tool integrates multiple predictive models with fuzzy context modeling to support case prioritization in domains where machine learning predictions must be complemented by expert-defined contextual criteria. CADEMAS-ML assigns weights to predictive models according to selected performance indicators and aggregates their outputs into a global risk score. In parallel, the tool evaluates contextual suitability using configurable fuzzy membership functions, categorical mappings, derived rules, and aggregation operators. This approach enables decision-makers to represent vague, gradual, and scenario-dependent preferences in a flexible manner. The final prioritization score combines the predictive risk component and the contextual suitability component through an adjustable trade-off parameter. The capabilities of the tool are illustrated through an employee attrition classification case study. Several departmental predictive models are combined with alternative decision contexts, including economic crisis and digital transformation scenarios. The results show that identical predictive risks may lead to different prioritization outcomes depending on the contextual assumptions adopted. This behavior improves the transparency and adaptability of cooperative decision-making processes. CADEMAS-ML provides a practical framework that connects predictive analytics, fuzzy reasoning, and human-configurable decision policies for the development of advanced decision-support systems.}

% Keywords
\keyword{automated decision-making systems; fuzzy decision making; context-aware decision support; machine learning; prioritization} 

%%%%%%%%%%%%%%%%%%%%%%%%%%%%%%%%%%%%%%%%%%
\begin{document}

\section{Introduction}\label{sec:intro}

Automated decision-making (ADM) systems based on machine learning (ML) have become an integral part of modern organizations. Predictive and classification models are routinely used to support tasks such as credit scoring, fraud detection, medical triage, customer churn analysis, and human-resource management. Their adoption has increased substantially in recent years due to their ability to process large volumes of data and identify patterns that support operational and strategic decisions.

Despite these advances, an important limitation remains. The outputs generated by ML models, such as class probabilities or risk scores, are often interpreted in isolation from the organizational environment in which decisions are ultimately made \citep{Tariq2024,Varga2023}. In practice, a high probability of an adverse event rarely provides sufficient guidance for action. The appropriate response depends on contextual factors such as organizational priorities, resource constraints, departmental objectives, and prevailing strategic conditions. As a result, identical predictive outcomes may lead to different decisions under different circumstances.

This observation highlights the cooperative nature of real-world decision-making \citep{Tariq2024}. Decisions are seldom based on a single predictive model. Organizations frequently rely on multiple models developed for different departments, objectives, or stakeholder groups. Their outputs must therefore be integrated into a coherent and comparable assessment of risk. At the same time, domain experts contribute contextual knowledge that cannot always be represented through precise numerical rules. Such knowledge is often expressed through gradual concepts and linguistic assessments that reflect uncertainty, preference, and strategic intent.

Fuzzy logic provides a well-established mathematical framework for representing this type of contextual reasoning \citep{Kumar2025,Ghosh2020}. Through fuzzy sets and linguistic variables, experts can model concepts such as urgency, strategic relevance, or suitability in a transparent and interpretable manner. This makes fuzzy approaches particularly attractive for decision-support environments in which quantitative predictions must be complemented by qualitative contextual considerations.

To address these challenges, \citet{NovoaHernandez2026} proposed \textbf{CADEMAS} (Cooperative Automated Decision-Making Systems), a conceptual framework that separates the decision process into three complementary components. The \emph{Input Manager} is responsible for collecting cases and predictive outputs. The \emph{Output Manager} aggregates the predictive information into a unified risk assessment. The \emph{Context Model} evaluates each case from a contextual perspective using fuzzy reasoning. Together, these components provide a structured mechanism for combining predictive evidence and contextual knowledge within a single decision-support process.

The practical value of this approach was subsequently illustrated by \citet{Godz2026}, who applied CADEMAS to employee-attrition prevention. Their study showed how predictive models developed by different organizational units could be coordinated with context-dependent management policies to generate prioritizations aligned with organizational objectives.

Although CADEMAS provides a solid conceptual foundation, its adoption is currently limited by the lack of a dedicated software implementation. Researchers and practitioners who wish to apply the framework must develop their own solutions to integrate predictive models, aggregate risk estimates, and construct fuzzy contextual assessments. Existing fuzzy and multi-criteria decision-making tools (e.g., \citealp{Wieckowski2022,Wieckowski2023,Roszkowska2024,Grazioso2023,Dikopoulou2021,Chen2020}) offer valuable functionalities, but none has been designed as an end-to-end cooperative ADM platform that combines multiple ML classifiers with a configurable fuzzy context layer. As discussed in Section~\ref{sec:related}, this creates a gap between the conceptual framework and its practical application.

This paper addresses that gap by presenting \textbf{CADEMAS-ML}, a software tool that operationalizes the CADEMAS framework for cooperative and context-aware ADM. The main contributions of this work are the following:

\begin{itemize}

\item Design and implementation of CADEMAS-ML, the first software tool that operationalizes the CADEMAS framework for cooperative and context-aware automated decision-making.

\item A configurable methodology for combining multiple machine-learning models and fuzzy contextual knowledge within a unified prioritization process.

\item A transparent and auditable decision-support architecture that exposes both predictive and contextual contributions to the final prioritization score.

\item An interactive robustness analysis mechanism that evaluates ranking stability under different balances between predictive evidence and contextual policies.

\item Validation through an employee-attrition prevention scenario involving multiple predictive models and alternative organizational contexts.

\end{itemize}

The remainder of the paper is organized as follows. Section~\ref{sec:background} reviews the CADEMAS framework and discusses recent related work on fuzzy and cooperative decision-support tools, with attention to whether they are released as open-source software. Section~\ref{sec:tool} describes CADEMAS-ML in detail, including its architecture, the predictive risk aggregation, the fuzzy contextual suitability layer, and the trade-off parameter. Section~\ref{sec:case} reports an employee-attrition case study under two contrasting contexts. Finally, Section~\ref{sec:conclusions} concludes and outlines directions for future work.

\section{Background and related works}

\subsection{CADEMAS framework}

The CADEMAS (Cooperative Automated Decision-Making Systems) framework was introduced as a general structure for designing decision-support systems that integrate multiple automated assessments while explicitly accounting for the organizational context in which decisions are made \citep{NovoaHernandez2026}. The central motivation behind CADEMAS is that real-world decision problems, especially in organizational environments, are rarely driven by a single predictive model. Instead, decisions typically emerge from the combination of multiple partial and sometimes conflicting evaluations produced by different units, models, or data sources, each with its own assumptions, information access, and reliability level.

To formalize this idea, let $\mathcal{E} = \{e_1, \dots, e_n\}$ denote a set of cases to be evaluated, and let $\mathcal{A} = \{A_1, \dots, A_K\}$ represent a collection of automated decision components. Each component produces a partial assessment $p_k(e_i) \in [0,1]$ for a given case $e_i$, which can be interpreted as a risk or priority indicator. Since these components may differ in performance and trustworthiness, each $A_k$ is assigned a confidence weight $w_k \geq 0$, with $\sum_{k=1}^{K} w_k = 1$. These weights can be derived from validation metrics, expert judgment, or governance constraints. The framework aggregates individual outputs into a global predictive risk score:
\begin{equation}
R(e_i) = \sum_{k=1}^{K} w_k \, p_k(e_i).    
\end{equation}

Rather than treating this aggregated value as the final decision criterion, CADEMAS introduces an explicit contextual layer that captures organizational priorities, policies, and scenario-dependent considerations. This layer is defined through a set of contextual factors $\mathcal{C} = \{C_1, \dots, C_L\}$, each of which encodes a relevant dimension of decision-making beyond predictive performance. These factors are modeled using fuzzy logic, allowing gradual degrees of satisfaction instead of binary rules. For each case $e_i$, the degree to which it satisfies factor $C_\ell$ is given by $\mu_\ell(e_i) \in [0,1]$, and these values are aggregated into a contextual score:
\begin{equation}
C(e_i) = H\big(\mu_1(e_i), \dots, \mu_L(e_i)\big),    
\end{equation}
where $H: [0,1]^L \rightarrow [0,1]$ is an aggregation operator that can take different forms depending on the decision policy (e.g., t-norms, averages, or hierarchical compositions).

The final prioritization index is obtained by combining predictive risk and contextual suitability:
\begin{equation}
\pi(e_i) = \Psi\big(R(e_i), C(e_i)\big),    
\end{equation}
where $\Psi$ is a monotonic function that defines the trade-off between both components. A key design principle of CADEMAS is the separation between predictive evidence and contextual interpretation, which makes explicit what the models suggest and what the organization considers relevant.

The framework has been instantiated in different application domains, including employee attrition prediction as described by \citet{Godz2026}. In that case, multiple machine learning models were trained on different organizational perspectives (e.g., Human Resources, Finance, Operations, and Employee Wellbeing). Each model acts as a decision component within $\mathcal{A}$, and their outputs are combined using weights derived from validation performance, specifically AUC. In parallel, the contextual layer encodes scenario-dependent organizational priorities (e.g. economic crisis or digital transformation settings) using fuzzy membership functions and rule-based structures.

Although this implementation demonstrates the practicality of the framework, it represents only one possible configuration. In general, CADEMAS is designed to remain flexible across several dimensions. Confidence weights $w_k$ can be defined using alternative performance measures (e.g., F1-score, precision, recall, calibration metrics, or utility-based indicators), depending on the operational objectives of the organization. Similarly, the aggregation operator $H$ is not restricted to a specific form and may be adapted to reflect different decision policies, ranging from conservative (conjunctive) to compensatory behaviors. The same applies to the combination function $\Psi$, which can be tuned to reflect different trade-offs between predictive and contextual information. Finally, contextual factors $\mu_\ell$ can be specified using different membership function families (triangular, trapezoidal, linear, or categorical) and may include conditional or scenario-dependent rules.

This flexibility is essential for adapting CADEMAS to heterogeneous organizational environments and motivates the need for a software implementation that operationalizes these design choices in a systematic and configurable way.

\subsection{Cooperative and fuzzy decision-support tools}

Recent research has increasingly recognized that effective decision support requires more than accurate predictions. In many application domains, decisions emerge from the interaction between predictive models, expert knowledge, and contextual considerations that may change over time. This observation has motivated the development of cooperative and human-centered decision-making frameworks that seek to combine the strengths of automated systems and human expertise.

Among these efforts, \citet{Tariq2024} propose an AI-to-collaborator (A2C) framework in which humans and AI systems share decision authority through explicit interaction mechanisms. Similarly, \citet{Varga2023} discuss cooperative AI architectures for high-stakes environments and argue that effective collaboration requires a shared representation of context between automated models and domain experts. These contributions provide important conceptual foundations for cooperative decision-making. However, they remain largely architectural or theoretical in nature and do not provide reusable software tools that practitioners can directly adopt. The CADEMAS framework proposed by \citet{NovoaHernandez2026} advances this line of research by explicitly separating predictive evidence from contextual reasoning through distinct architectural components.

At the same time, substantial progress has been made in the development of software tools for fuzzy decision support and multi-criteria decision making. Applications such as \textsc{MakeDecision} \citep{Wieckowski2024} and the fuzzy DEMATEL-based DSS described by \citet{Chekry2024} emphasize transparency and traceability by exposing intermediate calculations to users. In parallel, several open-source libraries have expanded the methodological toolbox available to researchers and practitioners. Examples include \textsc{pyFDM} \citep{Wieckowski2022}, \textsc{pyIFDM} \citep{Wieckowski2023}, \texttt{FuzzyR} \citep{Chen2020}, the intuitionistic-fuzzy toolkit of \citet{Roszkowska2024}, the \textsc{Eligere} platform \citep{Grazioso2023}, and the Fuzzy Cognitive Maps package presented by \citet{Dikopoulou2021}. These contributions have significantly improved the accessibility and reproducibility of fuzzy and MCDM methodologies. Nevertheless, their primary focus remains the implementation of fuzzy reasoning, ranking, or preference modelling techniques. They are not designed to integrate multiple predictive models within a unified decision-support workflow.

A related stream of research has explored the combination of machine learning and fuzzy reasoning. The survey by \citet{Kumar2025} highlights the growing interest in hybrid fuzzy-ML systems and identifies configurability and expert participation as key challenges for future decision-support tools. Other studies have proposed application-specific integrations of predictive models and expert knowledge, including fuzzy decision-support architectures \citep{Ghosh2020}, Fermatean fuzzy decision frameworks \citep{Suvitha2025}, and context-sensitive prioritization approaches based on fuzzy AHP and machine learning \citep{Aly2018}. Collectively, these works demonstrate the value of combining data-driven prediction with human knowledge. However, most of the reported solutions are tailored to particular domains and are presented as bespoke implementations rather than configurable and reusable software platforms.

\begin{table}[t]
\caption{Comparison of representative recent works on fuzzy and cooperative decision-support tools. The Software column reflects the availability of the associated implementation, when reported by the authors.\label{tab:related}}
\small
\begin{adjustwidth}{-\extralength}{0cm}
	\begin{tabularx}{\fulllength}{>{\raggedright\arraybackslash}l >{\raggedright\arraybackslash}X >{\raggedright\arraybackslash}X >{\raggedright\arraybackslash} l}
		\toprule
		\textbf{Work} & \textbf{Approach} & \textbf{Methodology} & \textbf{Software} \\
		\midrule
		\citet{NovoaHernandez2026} & Cooperative ADM framework & Coop., context-aware, fuzzy & Conceptual \\
		\citet{Godz2026} & CADEMAS in HR attrition & Coop., fuzzy, ML & Application \\
		\citet{Tariq2024} & Human--AI coop.\ (A2C) & Coop., hybrid & n/a \\
		\citet{Varga2023} & Cooperative AI & Coop., conceptual & n/a \\
		\citet{Wieckowski2024} & Fuzzy MCDM web app & Fuzzy MCDM, auditable & Web app \\
		\citet{Chekry2024} & Fuzzy DEMATEL DSS & Fuzzy MCDM & n/a \\
		\citet{Wieckowski2022} & pyFDM library & Fuzzy MCDM & OSS (Python) \\
		\citet{Wieckowski2023} & pyIFDM library & IF fuzzy MCDM & OSS (Python) \\
		\citet{Roszkowska2024} & IF-MCDM toolkit & Fuzzy MCDM & OSS \\
		\citet{Grazioso2023} & Eligere platform & Fuzzy ranking & OSS \\
		\citet{Dikopoulou2021} & FCM package & FCM, context model. & OSS \\
		\citet{Chen2020} & FuzzyR (R) & Fuzzy inference & OSS (R) \\
		\citet{Kumar2025} & Fuzzy--ML survey & Fuzzy, ML, hybrid & Survey only \\
		\citet{Ghosh2020} & Fuzzy DSS w/ rules & Fuzzy, hybrid & n/a \\
		\citet{Suvitha2025} & Fermatean fuzzy ops & Fuzzy aggregation & n/a \\
		\citet{Aly2018} & Fuzzy AHP + ML & Fuzzy MCDM, ML & n/a \\
		\textbf{CADEMAS-ML (this work)} & Cooperative ADM tool & Coop., fuzzy, ML, context-aware & OSS planned (Py.) \\
		\bottomrule
	\end{tabularx}
\end{adjustwidth}
\end{table}

Table~\ref{tab:related} summarizes representative contributions in these areas. Two observations emerge from this comparison. First, the software ecosystem surrounding fuzzy decision support has matured considerably, with many authors releasing reusable implementations and open-source libraries. Second, despite the methodological richness of the literature, no reviewed work simultaneously provides a cooperative decision-making architecture, the aggregation of multiple ML predictive models, a configurable fuzzy context layer, and an operational software environment that allows users to control the balance between predictive and contextual evidence.

CADEMAS-ML is designed to address this gap. Building upon the conceptual framework introduced by \citet{NovoaHernandez2026} and the employee-attrition application reported by \citet{Godz2026}, it provides a configurable software implementation in which predictive risk assessment and contextual reasoning coexist as explicit and auditable components of the decision-making process.


\section{CADEMAS-ML software tool}

The CADEMAS framework, as presented in Section~\ref{sec:cademas}, provides a formal and flexible architecture for cooperative, context-aware decision-making. However, its practical adoption requires a concrete implementation that streamlines the design, execution, and analysis of such systems across diverse application domains. This section introduces CADEMAS-ML, a software tool that translates the conceptual CADEMAS architecture into a functional, interactive web application. CADEMAS-ML operationalizes the three core layers—cooperative risk aggregation, fuzzy context evaluation, and decision integration—while explicitly supporting the flexibility discussed in Section~\ref{sec:cademas}: configurable performance metrics for confidence weighting, declarative fuzzy rule definitions, multiple aggregation operators, and adjustable trade-off policies. The tool is designed to empower analysts and decision-makers to instantiate, compare, and audit context-aware prioritization pipelines without requiring low-level programming. The following subsections detail the architectural design and the computational components that constitute the tool. A companion illustrative case study is presented in Section~\ref{sec:casestudy} to demonstrate its capabilities in a concrete organizational setting.

\subsection{Architecture}

CADEMAS-ML follows a modular, end-to-end pipeline architecture that separates data ingestion, model inference, context evaluation, and decision integration into independent but composable stages (Figure~\ref{fig:cademas_arch}). The tool is implemented as a web-based application using the Streamlit framework\footnote{\url{https://streamlit.io/}}, which enables interactive exploration and real-time parameter tuning without requiring local software installation. The source code is publicly available in an open-source repository\footnote{\url{https://github.com/...}}, and a live deployment is accessible online\footnote{\url{https://...}}.

The pipeline accepts four mandatory inputs: (i) a case dataset in CSV format (e.g., employee records); (ii) a collection of H2O MOJO models\footnote{\url{https://h2o.ai/}} representing the automated decision components; (iii) a JSON configuration specifying the performance metrics used to compute confidence weights; and (iv) a JSON configuration defining the fuzzy context rules. Upon execution, the pipeline performs batch inference over all cases, computes the cooperative risk score $R_i$ and the context alignment score $C_i$ for each case, and combines them into the final prioritization index $\pi_i$. All intermediate artifacts—per-model probabilities, membership degrees, and component contributions—are retained in session state for interactive inspection across multiple visualization tabs. This design ensures full traceability of the decision process while enabling comparative analysis under alternative configurations.

\begin{figure}[t]
\centering
% \includegraphics[width=0.9\textwidth]{figures/cademas_architecture.png}
\fbox{%
\parbox{0.8\textwidth}{\centering \vspace{3cm} 
\textbf{Figure placeholder:} 
Architectural pipeline of CADEMAS-ML showing data ingestion, MOJO ensemble inference, fuzzy context engine, $\lambda$-controlled integration, and visualization modules.%
\vspace{3cm}}
}
\caption{Modular architecture of CADEMAS-ML. The pipeline processes input data through parallel risk and context evaluation paths, combines them via the $\lambda$ trade-off parameter, and provides interactive visualization and robustness analysis.}
\label{fig:cademas_arch}
\end{figure}

\subsection{Predictive risk}

The cooperative risk layer implements the weighted ensemble aggregation defined in the CADEMAS framework. Let $\mathcal{M} = \{1, \dots, M\}$ denote the set of uploaded H2O MOJO models, each corresponding to an automated decision component $A_k$. For a user-selected performance metric $m$ (e.g., AUC, F1-score, log-loss), let $s_j$ be the performance value of model $j$ on metric $m$. The confidence weight for each component is computed by normalizing the performance scores:

\[
w_j = \frac{s_j}{\sum_{k \in \mathcal{M}} s_k}, \quad j \in \mathcal{M}.
\]

This weighting scheme reflects the relative trust assigned to each component based on prior validation evidence. The tool imposes no restriction on the number of models or the performance metric used, directly supporting the flexibility discussed in Section~\ref{sec:cademas}. For each case $i$, each model $j$ produces a positive-class probability $\hat{p}_{ij} \in [0,1]$ via MOJO inference. The global cooperative risk score is then computed as the weighted average:

\[
R_i = \sum_{j \in \mathcal{M}} w_j \, \hat{p}_{ij}.
\]

All models are required to agree on the predicted positive label across the entire batch; otherwise, the pipeline halts with a clear diagnostic error. Individual model probabilities are preserved in the session state, enabling per-model inspection in the dedicated \texttt{Models} tab. This transparency allows analysts to identify discrepancies among decision components and assess the influence of each model on the final risk score. The tool automatically handles feature alignment across heterogeneous models—a critical requirement when components operate on different feature subsets, as in the employee attrition scenario—by passing the full feature matrix to each MOJO model, leveraging H2O's internal alignment capabilities.

\subsection{Fuzzy context}

The context layer is the most distinctive and configurable component of CADEMAS-ML. It provides a declarative fuzzy inference engine that maps raw case features to a scalar context alignment score $C_i \in [0,1]$. This engine implements the fuzzy proposition model described in Section~\ref{sec:cademas} through a structured JSON configuration that can be defined and modified without code changes, embodying the flexibility of membership function families, aggregation operators, and conditional rules highlighted in the framework discussion.

The configuration schema comprises three hierarchical elements:

\begin{enumerate}
\item \textbf{Atomic rules}. Each atomic rule defines a mapping from a single feature to a membership degree $\mu_r(i) \in [0,1]$. The tool supports numeric membership functions (triangular, trapezoidal, linear increasing, linear decreasing), categorical direct mappings, and categorical set assignments. Importantly, atomic rules can include a \texttt{when} clause that restricts their activation to specific subpopulations (e.g., department-specific salary bands), enabling context-dependent policies without duplicating features.

\item \textbf{Derived rules}. Derived rules compose the outputs of one or more atomic or other derived rules through fuzzy aggregation operators. The tool implements four operators: \texttt{AND} (minimum, Gödel t-norm), \texttt{OR} (maximum, Gödel s-norm), \texttt{PRODUCT} (product t-norm), and \texttt{AVERAGE}. This allows analysts to construct hierarchical policy structures—for instance, defining \emph{commitment} as the minimum of job involvement and job satisfaction, and \emph{retention value} as the average of commitment and performance—directly in the configuration.

\item \textbf{Final aggregation}. The context score $C_i$ is obtained by aggregating the outputs of all atomic and derived rules. The tool provides two paths: (i) a user-selectable global operator (average, minimum, or product) applied across all numeric membership columns at runtime; or (ii) a declarative logic tree specified in the configuration, which defines a recursive expression over named rules. The former supports rapid interactive exploration; the latter ensures portability and reproducibility of context definitions across sessions.
\end{enumerate}

Every computed membership degree—atomic and derived—is stored in a per-case audit matrix, providing complete traceability of the fuzzy inference chain. The tool visualizes membership functions for numeric rules and presents the full audit table in the \texttt{Context} tab, enabling analysts to verify that the context configuration behaves as intended for individual cases. This auditability directly addresses the need for transparency in automated decision systems.

\subsection{Decision integration and robustness analysis}

The final prioritization index combines the cooperative risk and the context alignment score through a convex combination parameterized by a user-controlled trade-off factor $\lambda \in [0,1]$:

\[
P_i = \lambda \, R_i + (1 - \lambda) \, C_i,
\]

where $P_i$ denotes the prioritization score for case $i$. This formulation directly implements the generic combination function $\Psi$ from the CADEMAS framework (Section~\ref{sec:cademas}). The interpretation of $\lambda$ is intuitive: $\lambda = 1$ yields purely risk-driven prioritization, $\lambda = 0$ yields purely context-driven prioritization, and intermediate values represent explicit, adjustable trade-offs between predictive evidence and organizational policy. The tool visualizes the decomposition of $P_i$ into its ML contribution $\lambda R_i$ and context contribution $(1-\lambda) C_i$, making the influence of each component transparent at the individual case level.

Beyond static prioritization, CADEMAS-ML incorporates a dedicated \textbf{Robustness} module that assesses the stability of case rankings under varying decision policies. For a grid of $\lambda$ values, typically $\{0, \tfrac{1}{k}, \dots, 1\}$, the tool recomputes the prioritization scores and the induced ordering over cases. The results are presented as a bump chart, which visualizes how individual cases move up or down the priority list as the trade-off between risk and context shifts. This sensitivity analysis provides decision-makers with critical insight into whether their prioritization decisions are robust to reasonable variations in policy weighting, or conversely, whether certain cases are highly sensitive to the chosen $\lambda$ and thus require closer scrutiny. The module effectively operationalizes the flexibility of the $\Psi$ combination function, allowing stakeholders to explore the entire spectrum of possible decision policies within a unified interface.

\subsection{Implementation and availability}

CADEMAS-ML is implemented entirely in Python, leveraging the Streamlit framework for the interactive user interface and the H2O-3 Python API for MOJO model inference. The fuzzy context engine is implemented as a standalone module that parses JSON configurations and evaluates membership functions and aggregation operators in a vectorized manner over batches of cases. This design ensures that the tool remains responsive even with moderately large datasets (tested with thousands of cases and dozens of models). The application is containerized and can be executed locally or deployed on any cloud platform that supports Streamlit applications. The source code, example configurations, and a sample employee attrition dataset are provided in the public GitHub repository. A live demonstration instance is maintained to allow readers and practitioners to explore the tool's capabilities without installation. The tool is released under an open-source license to encourage adoption, extension, and integration with other decision-support ecosystems.

\section{A Case Study: Employee Attrition}

\section{Conclusion and future works}


\reftitle{References}

\bibliography{biblio.bib}


%\PublishersNote{}

\end{document}

